\documentclass{report}
\usepackage{amsmath,amssymb}
\setlength{\parindent}{0mm}
\setlength{\parskip}{1em}
\begin{document}
\begin{center}
\rule{6in}{1pt} \
{\large Meghan O'Connell \\
{\bf Krylov Recycling for Sequences of Shifted Systems}}

503 Boston Ave \\ Medford \\ MA \\ 02155
\\
{\tt meghan.oconnell@tufts.edu}\\
Misha Kilmer\end{center}

Solving many large scale sequences of shifted linear systems with
multiple right-hand sides is the main computational bottleneck of many
problems, such as nonlinear inverse problems. Specifically, this issue
arises in the optimization problem found in diffuse optical tomography
(DOT). An inner-outer Krylov recycling method was developed by Kilmer et
al. for the non-shifted case to reduce this computational cost. In this
paper, we extend this approach to the shifted case for a single
right-hand side. Both real shifts as well as complex identity shifts are
investigated. We show the value of our approach with two examples from
DOT, however, our method has the potential to be useful for other
applications as well.


\end{document}
